\section{Related Work}
\label{sec:related_work}

\subsection{Molecular Representation and Ligand Decomposition}

Molecular activity prediction models commonly employ two principal classes of
representations: fingerprint-based and graph-based representations. Morgan
fingerprints, commonly implemented as extended-connectivity fingerprints
(ECFPs), encode atomic neighborhoods as fixed-length bit vectors
\cite{rogers2010extended}. This representation is computationally inexpensive,
integrates well with classical machine-learning models, and remains a strong
baseline for many QSAR tasks. In contrast, graph neural networks (GNNs) learn
atomic and bond representations directly through message passing, thereby
reducing reliance on manually engineered features \cite{gilmer2017neural}.

In addition to whole-molecule representations, decomposing a ligand into
substructures helps relate model predictions to more readily interpretable
chemical motifs. BRICS cleaves molecules at bonds that can be synthetically
recombined according to chemical rules \cite{degen2008art}, whereas the
Bemis--Murcko scaffold retains ring frameworks and the linkers between rings
\cite{bemis1996properties}. Functional groups can be detected using SMARTS
queries. These three sources of substructures are complementary: BRICS describes
synthetically meaningful fragments, Murcko describes the core scaffold, and
SMARTS describes functional groups likely to govern chemical interactions.
Because substructures may overlap in their atom sets, they are treated as a set
of ligand ``views'' rather than as a single disjoint partition.

\subsection{Activity Regression and Pooling Mechanisms}

In classical QSAR, Ridge, Lasso, and Elastic Net are commonly applied to
fingerprints or descriptors. Regularization controls the high dimensionality and
collinearity of ECFPs, while linear coefficients provide a direct form of
interpretation. However, a linear model requires an operation that pools
substructure vectors into a ligand vector, such as the sum, mean, or maximum.

Attention models replace fixed pooling with learned weights. In attention
pooling, each substructure receives a weight, and the ligand representation is
the weighted sum of the embeddings. The Set Transformer extends this idea by
applying self-attention among the elements of an unordered set
\cite{lee2019set}. For molecular data, this approach enables interactions among
substructures to be modeled before constructing the whole-ligand representation.

\subsection{Substructure-Level Prediction Explanations}

Three families of explanation methods are investigated in the repository.
First, linear models decompose predictions according to regression coefficients
and feature vectors. Second, attention scores are used as proxies for
importance. Although convenient, this interpretation warrants caution because
attention does not necessarily correspond to causality or output sensitivity
\cite{jain2019attention}. Third, perturbation/masking directly measures the
change in prediction when each substructure is removed. Masking approximates the
counterfactual question, ``How would the prediction change if this motif were
absent?'', but the result may be affected by the molecule produced after atom
deletion, overlapping substructures, and samples outside the training domain.


\section{Methodology}
\label{sec:methodology}

\subsection{Problem Formulation and Data}

Given a ligand with SMILES \(x\), the objective is to learn the function
\(f_{\theta}(x)=\widehat{y}\), where \(y=\mathrm{pIC}_{50}\). The repository
uses the FAAH, CB2, PARP1, PBP2a, and USP7 datasets with predefined train/test
splits. Their sizes are summarized in Table~\ref{tab:datasets}. An augmented
version of FAAH was also created by generating 50 randomized SMILES for each
molecule, but no corresponding experimental log is available; this version is
therefore excluded from the quantitative comparison.

\begin{table}[htbp]
    \centering
    \caption{Datasets used in the study.}
    \label{tab:datasets}
    \begin{tabular}{lrrr}
        \hline
        Dataset & Train & Test & Total samples \\
        \hline
        FAAH   & 1075 & 120 & 1195 \\
        CB2    & 3511 & 391 & 3902 \\
        PARP1  & 2208 & 246 & 2454 \\
        PBP2a  & 80   & 9   & 89 \\
        USP7   & 458  & 51  & 509 \\
        \hline
    \end{tabular}
\end{table}

\subsection{Shared Ligand Decomposition}

Each SMILES is converted into an RDKit molecule and decomposed using three
mechanisms:
\begin{enumerate}
    \item \textbf{Murcko scaffold}: extract the core scaffold and map it back to
    the ligand atom indices;
    \item \textbf{BRICS}: identify bonds that can be cleaved according to BRICS
    rules, remove these bonds, and use the connected components as fragments;
    \item \textbf{Functional groups}: match 15 SMARTS patterns, including
    amides, esters, amines, and other common functional groups.
\end{enumerate}
Each substructure \(s_i\) is stored by its type, name, atom-index set, and the
canonical SMILES of its induced subgraph. Substructures with identical atom sets
are deduplicated, although substructures with different atom sets may still
overlap. If three-dimensional coordinates are required for the GNN, a conformer
is generated and optimized using MMFF when the input lacks one.

The repository also supports a legacy variant based on atom paths of lengths 2
to 6. However, the default experimental scripts all use the
BRICS--Murcko--functional-group combination; consequently, the results in
Section~\ref{sec:results} do not reflect the path-fragment variant.

\subsection{Substructure Representation}

For fingerprint-based methods, each \(s_i\) is encoded as an ECFP
\(\mathbf{e}_i\). The default configuration uses radius 2 and 2048 bits; the
baseline grid search examines radii \(\{2,3\}\) and bit counts
\(\{2048,4096\}\).

For the graph-based method, the entire ligand is passed through a GNN pretrained
on QM9. The GNN is frozen during the pIC50 experiments. The substructure
embedding is the mean of the embeddings of the atoms belonging to that
substructure:
\begin{equation}
    \mathbf{e}_i =
    \frac{1}{|A_i|}\sum_{v\in A_i}\mathbf{h}_v,
\end{equation}
where \(A_i\) is the atom set of \(s_i\), and \(\mathbf{h}_v\) is the atomic
embedding produced by the GNN. Batches containing different numbers of
substructures are padded and accompanied by a Boolean mask that excludes the
padding.

\subsection{pIC50 Prediction Models}

\paragraph{M1: Linear baseline.}
Substructure ECFPs are pooled using the sum, mean, and/or maximum. The full
configuration produces the vector
\begin{equation}
    \mathbf{z} =
    \left[
    \sum_i\mathbf{e}_i;\
    \frac{1}{m}\sum_i\mathbf{e}_i;\
    \max_i\mathbf{e}_i
    \right].
\end{equation}
After StandardScaler transformation, \(\mathbf{z}\) is passed to Ridge, Lasso,
or Elastic Net. The grid search comprises \(3\) regressors, \(2\) pooling
schemes, \(2\) radii, and \(2\) fingerprint sizes. The configuration is selected
using the mean \(R^2\) from 5-fold cross-validation on the training set.

\paragraph{M2: ECFP + attention pooling.}
Each ECFP is projected through a hidden layer:
\begin{align}
    \mathbf{h}_i &= \tanh(\mathbf{W}_h\mathbf{e}_i+\mathbf{b}_h),\\
    \alpha_i &=
    \frac{\exp(\mathbf{w}_a^\top\mathbf{h}_i)}
    {\sum_{j=1}^{m}\exp(\mathbf{w}_a^\top\mathbf{h}_j)},\\
    \mathbf{z} &= \sum_{i=1}^{m}\alpha_i\mathbf{h}_i.
\end{align}
A regression MLP maps \(\mathbf{z}\) to pIC50.

\paragraph{M3: GNN + substructure attention.}
M3 uses the same attention head as M2 but replaces ECFPs with substructure
embeddings obtained by mean-pooling the outputs of a pretrained, frozen atomic
GNN. This design tests whether graph-learned representations outperform fixed
fingerprints when the substructure pooling mechanism is held comparable.

\paragraph{M4: Set Transformer.}
The set of ECFP embeddings is processed by two Transformer encoder layers. A
learned query applies multi-head attention to the contextualized representations
to construct a ligand vector, after which an MLP predicts pIC50. The attention
weights returned by the pooling layer are averaged across heads by the
MultiheadAttention implementation.

\paragraph{M5: Masking attribution.}
M5 is not an independent prediction model and has no trainable parameters. This
method wraps a trained predictor, sequentially deletes the atoms of each
substructure, sanitizes the remaining molecule, and predicts pIC50 again.

\subsection{Training, Postprocessing, and Evaluation}

M1 standardizes features using StandardScaler within the pipeline. For M2--M4,
the means and standard deviations of the features and pIC50 are estimated from
the training set only. The models are optimized using MSE on the standardized
target; predictions are then transformed back to the pIC50 scale:
\begin{equation}
    \widehat{y} =
    \widehat{y}_{\mathrm{norm}}\sigma_{y,\mathrm{train}}
    +\mu_{y,\mathrm{train}}.
\end{equation}
M2 and M4 use 1000 epochs, a batch size of 128, and a learning rate of
\(10^{-4}\); M3 uses 200 epochs, a batch size of 32, and a learning rate of
\(10^{-3}\). The default seed is 42.

The primary metric is
\begin{equation}
    R^2=1-\frac{\sum_n(y_n-\widehat{y}_n)^2}
    {\sum_n(y_n-\bar{y})^2}.
\end{equation}
Computing \(R^2\) on the standardized target or on the original scale yields the
same value when both the target and predictions undergo a consistent affine
transformation. However, M2--M4 save the checkpoint with the best \(R^2\)
directly on the test set, whereas M1 selects its configuration by
cross-validation on the training set. The DL results may therefore be
optimistically biased because the test set participates in checkpoint selection
and does not constitute a fully independent evaluation.

\subsection{Post-Training Importance Estimation}

\paragraph{M1: Linear contribution.}
Because the model is fitted after StandardScaler transformation, the effective
coefficient for the original features is
\(\mathbf{w}_{\mathrm{eff}}=\mathbf{w}/\boldsymbol{\sigma}\). The contribution
of each substructure is computed separately for each pooling block and then
summed. For sum pooling,
\begin{equation}
    I_i^{\mathrm{sum}} =
    \mathbf{w}_{\mathrm{eff}}^\top\mathbf{e}_i.
\end{equation}
Mean pooling introduces an additional factor of \(1/m\). For max pooling, the
contribution in each dimension is assigned to the substructure attaining the
maximum according to the code's tie-handling rule. This implemented calculation
is more detailed than the simplified formula
\(I_i=\mathbf{w}^\top\mathbf{e}_i\) in the README.

\paragraph{M2 and M3: Attention score.}
The current inference path directly uses
\begin{equation}
    I_i=\alpha_i.
\end{equation}
Although the repository includes a helper that computes
\(\alpha_i\lVert\mathbf{e}_i\rVert_2\), and the README describes this formula
for M2, that helper is not invoked in the inference pipeline. The reported
results must therefore be interpreted as attention weights, rather than
attention multiplied by the embedding norm.

\paragraph{M4: Pooling attention.}
Importance is the attention weight from the learned query to the substructure:
\begin{equation}
    I_i=\operatorname{AvgHead}\left(a_{hi}\right).
\end{equation}
The implementation returns weights already averaged across heads by
MultiheadAttention and does not retain separate attributions for individual
heads.

\paragraph{M5: masking.}
For a trained predictor \(f\), signed importance is defined as
\begin{equation}
    I_i =
    f(x)-f(x\setminus s_i)
    =\widehat{\mathrm{pIC}}_{50}^{\,\mathrm{full}}
    -\widehat{\mathrm{pIC}}_{50}^{\,(-i)}.
\end{equation}
\(I_i>0\) indicates that deleting \(s_i\) decreases the predicted pIC50, whereas
\(I_i<0\) indicates that deleting the substructure increases the prediction.
For atom-level visualization, the repository distributes \(I_i\) equally among
the atoms in \(s_i\) and then sums contributions from overlapping substructures.


\section{Results and Discussion}
\label{sec:results}

\subsection{pIC50 Prediction Performance}

Table~\ref{tab:r2_results} summarizes the best test \(R^2\) values recorded in
\texttt{grid\_results.csv} and the \texttt{training.log} files. For M1, the
reported value corresponds to the configuration selected by cross-validation;
for M2--M4, it corresponds to the checkpoint with the best test \(R^2\).

\begin{table}[htbp]
    \centering
    \caption{Test \(R^2\) of the models on five datasets. Bold values indicate
    the highest score for each dataset.}
    \label{tab:r2_results}
    \begin{tabular}{lrrrr}
        \hline
        Dataset & M1: Linear & M2: FP-Attn & M3: GNN-Attn
                     & M4: Set Transformer \\
        \hline
        FAAH  & 0.414 & 0.453 & 0.097 & \textbf{0.461} \\
        CB2   & 0.469 & \textbf{0.535} & 0.056 & 0.533 \\
        PARP1 & $\approx 0.000$ & 0.002 & \textbf{0.005} & -0.009 \\
        PBP2a & 0.669 & \textbf{0.850} & 0.666 & 0.408 \\
        USP7  & \textbf{0.536} & 0.474 & 0.249 & 0.225 \\
        \hline
    \end{tabular}
\end{table}

No single architecture dominates across all datasets. FP-attention performs best
on CB2 and PBP2a and is close to the best result on FAAH. The Set Transformer
only marginally outperforms FP-attention on FAAH (0.461 versus 0.453), while
performing substantially worse on USP7 and PBP2a. The linear baseline remains
the best model on USP7, indicating that greater model complexity does not
necessarily improve performance when data are limited.

GNN-attention performs markedly worse than ECFP-attention on FAAH and CB2. One
plausible explanation is that the GNN was pretrained for QM9 property regression
and subsequently frozen, preventing its embeddings from adapting directly to
pIC50 or to the chemical domain of each target. These results are insufficient
to conclude that GNNs are generally inferior to fingerprints; they demonstrate
only that the current pretrained and frozen configuration is ineffective.

All models achieve values near \(R^2=0\) on PARP1, indicating no substantial
improvement over mean prediction on the current split. In contrast,
FP-attention achieves 0.850 on PBP2a. However, the PBP2a test set contains only
9 samples, so \(R^2\) has high variance and is highly sensitive to individual
observations. Furthermore, the M1 grid contains a configuration that achieves a
test \(R^2=0.848\), but this is not the configuration selected by CV; reporting
that value as the primary result would introduce test-based model-selection
bias.

\subsection{Qualitative Attribution Results}

The target-validation pipeline applies the models to ligands with PDB structures
and outputs the predicted pIC50 together with the substructures having the
largest \(|I_i|\). For example, on FAAH, the linear baseline predicts 8.22 for
ligand 2WAP and 8.84 for 3PPM, whereas the attention predictor used for masking
predicts 8.43 and 8.95, respectively. Under masking, one BRICS fragment of 2WAP
has \(\Delta\mathrm{pIC}_{50}=1.31\), and one BRICS fragment of 3PPM has
\(\Delta\mathrm{pIC}_{50}=1.68\).

Substructure rankings derived from linear contributions, attention, and masking
do not coincide completely. This is expected because the three methods answer
different questions: contribution within the fitted linear function, the weight
used when pooling embeddings, and the change in output following perturbation.
The repository currently provides neither ground-truth attributions nor a
quantitative measure of agreement among the explanation methods; the heatmaps
should therefore be regarded only as qualitative analyses.

\subsection{Experimental Limitations}

The preceding results should be interpreted in light of four limitations.
First, model checkpoints are absent from the current repository snapshot,
although the logs and result CSV files remain; inference cannot yet be
reproduced solely from the stored artifacts. Second, M2--M4 use the test set to
select the best epoch, compromising the independence of the evaluation. Third,
FAAH augmentation has been prepared, but no training results are available.
Finally, the repository does not report MAE, RMSE, confidence intervals across
multiple seeds, or external validation against observed pIC50 values. Future
experiments should separate the validation set from the test set, run multiple
seeds, preserve molecule-level splits when using randomized-SMILES augmentation,
and evaluate attribution using both masking consistency and independent
medicinal-chemistry knowledge.

\bibliographystyle{plain}
\bibliography{references}
